\chapter*{Abstract}
\addcontentsline{toc}{chapter}{Abstract}

\begin{spacing}{1.4}
This thesis presents an end-to-end American Sign Language (ASL)
recognition system that runs entirely on a consumer Android phone.
Two complementary tasks are addressed: fingerspelling recognition
over the 29 classes of the ASL Alphabet dataset, and word-level
recognition over the 100-class WLASL-100 vocabulary. For the letter
task, two classifiers are trained and compared: an EfficientNet-B0
that consumes raw image crops and a compact multi-layer perceptron
over 21 MediaPipe hand landmarks. On the single-signer ASL Alphabet
test split they reach \SI{100}{\percent} and \SI{99.38}{\percent}
Top-1 respectively, read as a pipeline sanity check rather than
capability claims; the landmark model ships in a \SI{0.24}{MB}
checkpoint. A single-user on-device manual evaluation under
realistic camera conditions inverts the offline ranking on the
budget phone (landmark MLP \SI{100}{\percent}, image network
\SI{92.9}{\percent}), motivating the landmark variant as the
deployment default for the letter task. For the word task, a
bidirectional LSTM is trained on a two-handed 128-dimensional
landmark feature stream extracted from a yt-dlp-recovered WLASL-100
corpus (\num{1465} clips, a \SI{44}{\percent} increase over the
\num{1013}-video Kaggle baseline). A five-seed logit-averaging
ensemble reaches \SI{56.22}{\percent} Top-1 and
\SI{81.09}{\percent} Top-5 on the WLASL-100 test partition under
the canonical person-mixed split, clearing the \SI{40}{\percent}
Top-1 acceptance bar adopted for this thesis by sixteen percentage
points. Because the protocol is person-mixed by design, this
offline figure is positioned as a literature-comparable benchmark
point under that protocol rather than as a signer-independent
generalization estimate. A nine-step forward-additive ablation
decomposes the gain from a \SI{20}{\percent} baseline to the
\SI{60.08}{\percent} final validation accuracy into three pillars
of comparable weight: data engineering, the extraction pipeline,
and training-recipe refinements. All four classifiers are exported
to TensorFlow Lite and deployed in an Android application built
around CameraX and MediaPipe Hands. On-device throughput, latency,
and the OOD accuracy gap are characterised on a budget
(TECNO SPARK 20C) and a flagship (Xiaomi 15T Pro); word accuracy
is flagship-only because entry-level ensemble latency is too high
for interactive use. The evaluation produces a device-conditional
model-selection guideline.
\end{spacing}
